% --- INICIO DE PLANTILLA PERSONALIZADA ---
\documentclass[oneside,11pt]{Latex/Classes/PhDthesisPSnPDF}

% --- Paquetes Originales ---
\usepackage{blindtext}
\usepackage{actuarialangle} 
\usepackage{tabularx}
\usepackage{float}
\usepackage{multirow, array}
\usepackage[usenames,dvipsnames,svgnames,table]{xcolor}
\usepackage{longtable}
\usepackage{latexsym,amsmath,amssymb,amsfonts}
\usepackage{enumitem}
\usepackage{xparse}
\usepackage{booktabs}
\usepackage[round, sort, numbers]{natbib}  
\usepackage{adjustbox}
\usepackage[utf8]{inputenc}
\usepackage{caption}
\usepackage{graphicx}
\usepackage{hyperref}
\usepackage{cleveref}

% Definición de colores
\definecolor{miblue}{RGB}{68, 114, 196}

% --- CAMBIO CRÍTICO 1: Usar \input para macros en el preámbulo ---
% \include da problemas antes del begin document. \input es seguro.
\input{Latex/Macros/MacroFile1}

% --- Definiciones de seguridad (por si la clase falla al cargar alguna variable) ---
\providecommand{\facultad}[1]{\def\lafacultad{#1}}
\providecommand{\escudofacultad}[1]{\def\elescudo{#1}}
\providecommand{\degree}[1]{\def\elgrado{#1}}
\providecommand{\director}[1]{\def\eldirector{#1}}
\providecommand{\lugar}[1]{\def\ellugar{#1}}
\providecommand{\degreedate}[1]{\def\lafecha{#1}}

% --- Configuración para bloques de código de R (Pandoc) ---

%%%%%%%%%%%%%%%%%%%%%%%%%%%%%%%%%%%%%%%%%%%%%%%%%%%%%%%%%%%%%%%%%%%%%%%%%%%%%%%%
%                         DATOS (Conectados a RMarkdown)                       %
%%%%%%%%%%%%%%%%%%%%%%%%%%%%%%%%%%%%%%%%%%%%%%%%%%%%%%%%%%%%%%%%%%%%%%%%%%%%%%%%
\title{Impacto de la adicción digital en el rendimiento académico y
bienestar de los estudiantes en diferentes países}
\author{Michelle Salazar \\ José Hernández \\ Diego Guerrero}

% Lógica para Facultad: Si la defines en el YAML la usa, si no, usa la default

\facultad{Facultad de Ciencias Económicas y Sociales \\ Escuela de
Estadistica y Ciencias Actuariales}
  \escudofacultad{Latex/Classes/Escudos/faces}

\degree{Computación I}
\director{Jesus Ochoa \\ Oliver Triveño}
\degreedate{Marzo 2026} 
\lugar{Caracas}

\portadatrue 

% Metadatos PDF
\keywords{}
\subject{}

% --- CORRECCIÓN PARA PANDOC NUEVO (Imágenes) ---
\makeatletter
\providecommand{\pandocbounded}[1]{#1}
\makeatother


%%%%%%%%%%%%%%%%%%%%%%%%%%%%%%%%%%%%%%%%%%%%%%%%%%%%%
%                   DOCUMENTO                       %
%%%%%%%%%%%%%%%%%%%%%%%%%%%%%%%%%%%%%%%%%%%%%%%%%%%%%
\begin{document}

\maketitle

%%%%%%%%%%%%%%%%%%%%%%%%%%%%%%%%%%%%%%%%%%%%%%%%%%%%%
%                  PRÓLOGO                          %
%%%%%%%%%%%%%%%%%%%%%%%%%%%%%%%%%%%%%%%%%%%%%%%%%%%%%
\frontmatter

% Puedes agregar dedicatorias desde el YAML si quieres, o dejarlas fijas aquí

%%%%%%%%%%%%%%%%%%%%%%%%%%%%%%%%%%%%%%%%%%%%%%%%%%%%%
%                   ÍNDICES                         %
%%%%%%%%%%%%%%%%%%%%%%%%%%%%%%%%%%%%%%%%%%%%%%%%%%%%%
\setcounter{secnumdepth}{3}
\setcounter{tocdepth}{3}

\tableofcontents



\cleardoublepage

  \cleardoublepage       % Empieza en página derecha (estándar en tesis)
  \chapter*{Resumen}     % Crea el título "Resumen" sin numeración (Capítulo *)
  \addcontentsline{toc}{chapter}{Resumen} % (Opcional) Lo añade al índice
  
  Resumen del trabajo de
investigación\ldots{}             % Aquí se pega el texto que escribiste en el YAML
  
  \cleardoublepage

%%%%%%%%%%%%%%%%%%%%%%%%%%%%%%%%%%%%%%%%%%%%%%%%%%%%%
%                   CONTENIDO (El Corazón)          %
%%%%%%%%%%%%%%%%%%%%%%%%%%%%%%%%%%%%%%%%%%%%%%%%%%%%%
\mainmatter
\def\baselinestretch{1.5}

% --- CAMBIO CRÍTICO 2: Aquí RMarkdown inyecta todo tu texto ---
\chapter*{Introducción}

Introducción del trabajo de investigación

\chapter{Planteamiento del Problema}

Describir planteamiento del problema~Desde su aparición las redes
sociales cambiaron la manera en que las personas interactúan con el
mundo real, desde servir como medios de comunicación hasta convertirse
en plataformas que afectan el comportamiento del ser humano. En los
estudiantes estas plataformas son útiles para acceder a información y
conocer noticias mundiales. Sin embargo, el uso de las redes sociales se
ha transformado en un fenómeno que crece a gran escala, la adicción a
las redes, y la adicción al mundo digital. El problema radica en que el
uso de las redes y plataformas digitales puede llegar a un extremo de
dependencia, e incluso excesiva en algunos casos. Siendo un factor que
influye en los estudiantes y en su capacidad de concentración, afectando
su rendimiento académico. La gran mayoría de estudiantes tienen acceso a
las diferentes plataformas de redes sociales, pero desconocemos su
comportamiento a nivel geográfico y variables como el tipo específico de
red social (TikTok, Instagram, WhatsApp, etc.) factores como la salud
mental, o el entorno geográfico puede darnos ciertas ``tendencias''
sobre la dependencia a las redes sociales. Por esto, es necesario
investigar si la adicción a las plataformas digitales y redes sociales
se comporta de la misma manera en todo el mundo, o responde a patrones
específicos por región.

Argumento final. Debido a esto, se plantea las siguientes preguntas de
investigación:\\

\textbf{1.} Pregunta1~¿Cómo se distribuyen los nivelees de adicción
digital con dependencia significativamente alta?

\textbf{2.} Pregunta2~¿Los estudiantes comn mayores puntajes de adicción
digital reportan mayor frecuencia que las redes solciales afectan
negativamente en su rendimiento académico?

\textbf{3.} Pregunta3~¿De que forma el uso de redes sociales especificas
afecta la salud mental y reduce las horas de sueño de los estudiantes?

\textbf{4.} Pregunta4~¿Cómo varían los niveles de adicción digital entre
los diferentes países y qué patrones de riesgo nos permiten clasificar a
los países según el impacto de ésta dependencia?

\section{Justificación}

Justificación del trabajo de investigación~El proposito ce esta
investigación es identificar los patrones y comportamientos del uso
excesivo de las redes sociales, identificando que niveles significativos
se encuentran y que relaciones tienen con la vida del estudiante,
comprender que riesgos surguen en su rendimiento académico, problemas de
salud mental y trastornos de sueño. Mediante éste análisis se busca
cuantificar dichas problemáticas que lleven a aplicar estrategias y
concientizar la prevención a la hora de usar redes sociales, logrando
establecer una vida más saludable, productiva y equilibrada para los
estudiantes en distintos países.

\section{Objetivos}

\subsection{Objetivo General}
\begin{itemize}
  \item{Analizar la adicción a las redes sociales y su impacto en el bienestar académico y mental de los estudiantes, comparando los patrones de comportamiento entre los diferentes países.}
\end{itemize}

\subsection{Objetivos Específicos}
\begin{itemize}
  \item{Describir la adicción estadística de los puntajes de adicción digital de los estudiantes mediante medidas de tendencia central y dispersión, para determinar si existen grupos con una dependencia significativamente alta}
  \item{Determinar la asociación entre el nivel de adicción digital y la percepción de como afecta el rendimiento academico para reconocer si los estuduantes con mayor dependencia reportan un imoacto negativo en sus estudios}
  \item{Determinar el impacto del uso de redes sociales en la salud mental y calidad de sueño de los estudiantes, identificando las redes sociales que afectan en mayor nivel}
  \item{Reconocer que plataformnas de redes sociales vincula con los niveles más bajos de salud mental y horas de sueño}
  \item{Identificar patrones geográficos de riesgo mediante la clasificación de países según sus niveles de adición permitiendo detectar tendencias regionales y variaciones significativas en el impacto de la dependencia digital. }
\end{itemize}

\section{Cobertura}

\subsection{Horizontal}

Describir la cobertura Horizontal del trabajo de investigación\\

\begin{table}[ht]
\centering
\caption{Variables consideradas} 
\resizebox{10cm}{!} {
\begin{tabular}{ccc}
  \hline
\ \ \ \ \ NOMBRE DE LA VARIABLE \\ 
  \hline
   VARIABLE 1 \\ 
VARIABLE 2\\ 
VARIABLE n\\ 
   \hline
\end{tabular}
}
\end{table}

\subsection{Vertical}

Describir la cobertura Vertical del trabajo de investigación

\section{Periodo de Referencia}

Describir la fuente y el Periodo de Referencia utilizado en el trabajo
de investigación

\section{Variables}

Describir la variable en estudio en el trabajo de investigación

\chapter{Marco Teórico}

\section{Bases Teóricas}

Las redes sociales son un medio que cumple la función de facilitar la
interacción de las personas. A través de las redes sociales, las
personas pueden cambiar ideas en tiempo real y compartir información en
cualquier parte del mundo. Aunque ya no son solamente un medio para
socializar, sino que estas redes tienen aplicaciones que afectan el
rendimiento académico y la vida cotidiana de los estudiantes. (Castells,
2012) En este sentido, el uso excesivo de redes sociales puede tener un
impacto considerable en el bienestar psicológico de los estudiantes
universitarios, afectando diversos aspectos de su salud mental. En éstas
plataformas, se observa un aumento de la ansiedad y la presión, sobre
todo cuando se prioriza la apariencia y la posición social. Los alumnos
a menudo sienten la necesidad de estar a la altura de sus compañeros, lo
que genera un ambiente de tensión y un mayor juicio entre ellos. Este
fenómeno es particularmente evidente en aquellos que buscan mantener una
reputación impecable o que compiten por estatus y reconocimiento social
(Campos, 2021). Así cómo también, dicho uso excesivo puede alterar los
hábitos de sueño. Un trabajo prolongado ante la pantalla y un uso
intenso de la red son factores que pueden alterar los ciclos de sueño de
quien está aprendiendo, un hecho que ocasiona que su desempeño académico
se vea afectado, así como su salud mental, un mal ciclo que puede llevar
a los estudiantes a la desesperanza (Campos, 2021) El rendimiento
académico de un alumno depende de múltiples variables, tanto internas
como externas; el impulso se encuentra entre factores internos, ya que
determina en que medida el alumno tiende a esforzarse de la manera que
requiere el aprendizaje. En ocasiones las personas quieren aprender por
su innata curiosidad, o solamente por impresionar. La consideración de
la motivación inherente y externa puede influenciar el rendimiento
educativo (Gonzalez, 2024). Por otra parte, se dice que lo que se conoce
como ``adicción a las redes sociales'' ha sido discutido por autores que
han tratado acerca de las ciberdirecciones, autores como Enrique
Echeburúa, Mark Griffiths, Xavier Sánchez-Carbonell, Rubicelia Valencia
Ortiz, etc. A finales de la década de los noventa se comenzó a hablar
sobre una posible adicción a la televisión y a los videojuegos, siendo
el doctor Mark Griffiths el primer autor en hablar sobre las
``adicciones tecnológicas'' (Griffiths, 1995). Luego se desarrollaron
ideas sobre los riesgos del uso del teléfono y el internet. Actualmente
con el auge de las redes sociales se habla sobre una posible
``adicción'' a las mismas. Analizando la perspectiva que defiende que el
trastorno a las redes sociales no es una ``adicción'' sino una necesidad
que lleva a un mal uso de manera excesiva de las redes sociales
(Sánchez-Carbonell, 2008), se dice que se basa en tres factores:

• En primer lugar, argumenta la diferencia de este tipo de trastornos
con respecto a las adicciones reconocidas por la OMS. Para realizar ésta
afirmación se basan en que no existen sustancias adictivas que
condicionen el funcionamiento psico-neuronal del usuario, estableciendo
que las malas conductas no son considerables como algo importante a la
hora de dictar si se considera una adicción o no. Un ejemplo de ello
sería el factor de horas invertidas en redes sociales, ¿podría
considerarse un factor fundamental?, la respuesta es no, debido a que el
usuario puede tener un uso excesivo de redes sociales sin ser
necesariamente adicto (Sánchez-Carbonell y Ursulla Obsert, 2015).

• En segundo lugar, no hay lugar a la comparación en el efecto directo
que poseen las drogas con sustancias, respecto a los comportamientos
adictivos. Las drogas influyen en la salud tanto mental como física del
usuario. Mientras que múltiples estudios establecen que la ``adicción a
las redes sociales'' no es causante de trastornos directamente, pero
repercute indirectamente en la creación de los mismos, como puede ser
depresión, miedo social ó impulsividad (Echeburúa, 2012).

• En tercer lugar, exponen la necesidad de adaptarse a una realidad
basada en la tecnología. Ésta va cambiando constantemente, la manera de
contactarse ha cambiado drásticamente y en la actualidad son más
constantes, pero a su vez menos intensas. Por lo que también está la
necesidad de adaptarnos y entender esa nueva forma de socializar.
Estableciendo que a lo mejor las personas no son adictas a las redes
sociales, si no que la sociedad actual demanda su uso constante, para
que el usuario no quede obsoleto o anticuado en la misma
(Sánchez-Carbonell y Ursulla Oberst, 2015).

Bajo éstos factores principales, se defiende la idea de que las redes
sociales no son adictivas, pero que no sean consideradas una adicción no
significa que no tengan en cuenta sus riesgos. Consideran que dichos
riesgos se pueden solucionar con el paso del tiempo, y la adaptación de
la sociedad a éstas tecnologías y formas de comunicación
(Sánchez-Carbonell y Ursulla Oberst, 2015). De igual manera, estos
autores defienden que las redes sociales no solamente tienen usos
perjudiciales para las personas que las consumen, sino que son de gran
utilidad al usarlas correctamente, pudiendo tener beneficios positivos
para su desarrollo social. Provocando el aumento de oportunidades
laborales y académicas, así como una mejor participación en la sociedad
(Kuss y Griffiths, 2011; Sánchez-Carbonell y Ursulla Oberst, 2015). El
concepto de adicción ha estado vinculado con las drogas, por eso este
término necesitaría una redefinición del mismo, adaptándolo a las
adicciones de comportamiento. Aunque no hay una definición concreta
sobre la adicción a las redes, tampoco hay unidad ya que los autores
realizan distintas conceptualizaciones del término. Según Echeburúa
(2012): Una adicción sin droga es toda aquella conducta repetitiva que
resulta placentera en las primeras fases, generando descontrol en el
sujeto, interrumpiendo su vida cotidiana a nivel social, laboral o
familiar (p.~436). Con las investigaciones de ambas posturas y la
definición de Echeburúa (que puede adaptarse a las redes sociales). Se
parte de la idea de que las adicciones de comportamiento tienen la misma
gravedad que las sustancias. Siendo el término correcto ``adicción a las
redes sociales'' y no un ``uso conflictivo''.

Resumen de las bases teoricas explicadas en el capitulo y su
importancia~

\subsection{Definición 1}

Descripción de la definición 1\\

\subsection{Definición 2}

Descripción de la definición 2\\

\subsection{Definición n}

Descripción de la definición n\\

\chapter{Marco Metódico}

En esta sección se presentan las metodologías y/o test necesarios
relacionados con la práctica divulgada en el marco teórico, así como las
fuentes de datos.

\section{Bases metódicas utilizadas en el trabajo de investigación}

Desarrollo\\

\subsection{Item 1}

Descripción\\

\subsection{Item 2}

Descripción\\

\subsubsection{Item 2.2}

Descripción\\

\subsection{Item n}

Descripción\\

\chapter{Análisis de Resultados}

\{\small Este capítulo presenta los resultados obtenidos al aplicar las
metodologías descritas en el capítulo anterior, así como del análisis
estadístico descriptivo de las variables estudiadas y las ventajas y
desventajas al aplicar dichos tests.\}

\section{Presentación de resultados 1}

Descripción\\

\section{Presentación de resultados 2}

Descripción\\

\section{Presentación de resultados n}

Descripción\\

\chapter{Conclusiones y Recomendaciones}

\textbf{1.-} Conclusion 1 del trabajo de investigación.\\

\textbf{Recomendación} Recomendación 1 del trabajo de investigación.\\

\textbf{2.-} Conclusion 2 del trabajo de investigación.\\

\textbf{Recomendación} Recomendación 2 del trabajo de investigación.\\

\textbf{n.-} Conclusion n del trabajo de investigación.\\

\textbf{Recomendación} Recomendación n del trabajo de investigación.\\

\appendix
\renewcommand{\thechapter}{\Alph{chapter}}

\chapter{Anexo A}

\chapter{Anexo B}

\chapter{Anexo C}

%%%%%%%%%%%%%%%%%%%%%%%%%%%%%%%%%%%%%%%%%%%%%%%%%%%%%
%                   REFERENCIAS                     %
%%%%%%%%%%%%%%%%%%%%%%%%%%%%%%%%%%%%%%%%%%%%%%%%%%%%%
% Mantenemos la estructura natbib de tu plantilla
\renewcommand{\bibname}{Lista de Referencias}
\bibliographystyle{apalike} 
\bibliography{bibliografia.bib} 

%%%%%%%%%%%%%%%%%%%%%%%%%%%%%%%%%%%%%%%%%%%%%%%%%%%%%
%                   APÉNDICES                       %
%%%%%%%%%%%%%%%%%%%%%%%%%%%%%%%%%%%%%%%%%%%%%%%%%%%%%

\end{document}
